\section{Conclusion}
We presented \algo, a simple, general interface that turns text-conditioned video predictions into executable robotic actions by reconstructing and tracking 3D object flow. This decouples what should happen in the world (task-relevant object motion and state change) from how a particular embodiment realizes it under kinematic, dynamic, and morphology constraints. Built entirely from off-the-shelf video generation and perception tools, \algo solves open-world manipulation tasks in simulation and on real robots across rigid, articulated, deformable, and granular objects using only RGB-D observations and language, without task-specific demonstrations. Experiments show consistent gains over trajectory baselines derived from dense optical flow or rigid pose transforms, robustness to variations in instances, backgrounds, and viewpoints, and the importance of both the upstream video model and downstream dynamics choice (with particle dynamics proving most reliable). A failure analysis highlights current bottlenecks, such as video artifacts (morphing, hallucinations), occlusion-induced tracking dropouts, and grasp selection mismatches as concrete directions for improvement.  Overall, our results indicate that 3D object flow is a scalable bridge from open-ended video generation to robot control in unstructured environments.
